\begin{center}\textit{Algorithmus novus ad decidendum, utrum numerus integer positivus datus numeri primi positivi dati residuum quadraticum sit an non-residuum.}\end{center}

\Article{1}

Antequam solutionem novam huius problematis exponamus, solutionem in \textit{Disquisitionibus Arithmeticis} traditam hic breviter repetemus, quae satis quidem expedite perficitur adiumento theorematis fundamentalis atque theorematum notorum sequentium:

I. Relatio numeri $a$ ad numerum $b$ (quatenus ille huius residuum qua\-draticum est sive non-residuum), eadem est quae numeri $c$ ad $b$. si $a\equiv c\ (\mathrm{mod.}\ b)$.

II. Si $a$ est productum e factoribus $\alpha,\beta,\gamma,\delta$ etc., atque $b$ numerus primus, relatio ipsius $a$ ad $b$ ita a relatione horum factorum ad $b$ pendebit, ut $a$ fiat residuum quadraticum ipsius $b$ vel non-residuum, prout inter illos factores reperitur multitudo par vel impar talium, qui sint non-residua ipsius $b$. Quoties itaque aliquis factor est quadratum, ad eum in hoc examine omnino non erit respiciendum; si quis vero factor est potestas integri cum exponente impari, illius vice ipse hic integer fungi poterit.

III. Numerus 2 est residuum quadraticum cuiusvis numeri primi formae $8m+1$ vel $8m+7$, non-residuum vero cuiusvis numeri primi formae $8m+3$ vel $8m+5$.

\SourceLeaf{70}%
Proposito itaque numero $a$, cuius relatio ad numerum primum $b$ quaeritur: pro $a$, si maior est quam $b$, ante omnia substituetur eius residuum minimum positivum secundum modulum $b$, quo residuo in factores suos primos resoluto, quaestio per theorema II reducta est ad inventionem relationis singulorum horum factorum ad $b$. Relatio factoris $2$, (siquidem adest vel semel, vel ter, vel quinquies etc.) innotescit per theorema III; relatio reliquorum, per theorema fundamentale, pendet a relatione ipsius $b$ ad singulos. Hoc itaque modo loco unius relationis numeri dati ad numerum primum $b$ iam investigandae sunt aliquae relationes numeri $b$ ad alios primos impares ipso $b$ minores, quae problemata eodem modo ad minores modulos deprimentur, manifestoque hae depressiones successivae tandem exhaustae erunt.

\Article{2}

Ut exemplo haec solutio illustretur, quaerenda sit relatio numeri $103$ ad $379$. Quum $103$ iam sit minor quam $379$, atque ipse numerus primus, protinus applicandum erit theorema fundamentale, quod docet, relationem quaesitam oppositam esse relationi numeri $379$ ad $103$. Haec iterum aequalis est relationi numeri $70$ ad $103$, quae ipsa pendet a relationibus numerorum $2,\ 5,\ 7$ ad $103$. Prima harum relationum e theoremate III innotescit. Secunda per theorema fundamentale pendet a relatione numeri $103$ ad $5$, cui per theorema I aequalis est relatio numeri $3$ ad $5$; haec iterum per theorema fundamentale pendet a relatione numeri $5$ ad $3$, cui per theorema I aequalis est relatio numeri $2$ ad $3$, per theorema III nota. Perinde relatio numeri $7$ ad $103$ per theorema fundamentale a relatione numeri $103$ ad $7$ pendet, quae per theorema I aequalis est relationi numeri $5$ ad $7$; haec iterum per theorema fundamentale pendet a relatione numeri $7$ ad $5$, cui aequalis est per theorema I relatio numeri $2$ ad $5$ per theorema III nota. Quodsi iam hanc analysin in synthesin transmutare placet, quaestionis decisio ad quatuordecim momenta refertur, quae complete hic apponimus, ut maior concinnitas solutionis novae eo clarius elucescat.

1.\quad Numerus $2$ est residuum quadraticum numeri $103$ (theor. III).

2.\quad Numerus $2$ est non-residuum quadraticum numeri $3$ (theor. III).

3.\quad Numerus $5$ est non-residuum quadraticum numeri $3$ (ex I et 2).

4.\quad Numerus $3$ est non-residuum quadraticum numeri $5$ (theor. fund. et 3).

5.\quad Numerus $103$ est non-residuum quadraticum numeri $5$ (I et 4).

